%% ============================================================
%%  生成AI利用申告書 — 発展課題 (HCK02_03.ino / HCK02_03.pde)
%%  ハッカソン1 / 塩澤 拓海 / 2026
%%  コンパイル: docker run --rm -v "$(pwd):/workspace" -w /workspace \
%%              ghcr.io/paperist/texlive-ja:debian latexmk main.tex
%% ============================================================
\documentclass[a4paper,11pt,dvipdfmx]{jsarticle}

\usepackage[top=25mm,bottom=25mm,left=25mm,right=25mm]{geometry}
\usepackage{amssymb}
\usepackage{listings}
\usepackage{plistings}
\usepackage[dvipdfmx,hidelinks]{hyperref}
\usepackage{pxjahyper}

\lstset{
  basicstyle=\ttfamily\small,
  frame=single,
  breaklines=true,
  xleftmargin=1em,
  aboveskip=6pt, belowskip=6pt,
  columns=flexible,
}

\newcommand{\cboxon}{$\boxtimes$}
\newcommand{\cboxoff}{$\square$}
\newcommand{\todo}[1]{\textbf{[TODO: #1]}}

\begin{document}

{\LARGE\bfseries 生成AI利用申告書}
\vspace{1em}

\begin{itemize}
  \item 氏名：塩澤 拓海
  \item 学籍番号：\todo{学籍番号を記入}
  \item プログラムファイル名：\texttt{HCK02\_03.ino} / \texttt{HCK02\_03.pde}
  \item 使用したAIツール：Claude Code (Opus 4.7)
\end{itemize}

\section*{1. 何に使ったか}
発展課題（マイク入力音信号を Processing でスペクトル解析し、結果をウィンドウに描画する
スペクトルアナライザを作る）の設計と実装で生成AIを利用した。提出課題2 (\texttt{HCK02\_02.ino}
/ \texttt{HCK02\_02.pde}) の構成を土台に、Processing 側を時間領域表示から
周波数領域 (FFT) 解析へ拡張する作業の補助として用いた。
\begin{itemize}
  \item Arduino 側 (\texttt{HCK02\_03.ino}) のサンプリング周期と Processing 側の
        \texttt{SAMPLE\_RATE\_HZ} を一致させるための数値設計
        （\texttt{SAMPLE\_US} = 150\,µs $\rightarrow$ 6666.67\,Hz、ナイキスト約 3333\,Hz）
  \item Processing 側 (\texttt{HCK02\_03.pde}) で minim の \texttt{ddf.minim.analysis.FFT}
        を使ってシリアル受信サンプルを FFT にかけ、ビン強度を dB 表示する設計の相談
  \item 時間領域波形と周波数領域スペクトルを 1 ウィンドウに上下 2 段で並べる
        レイアウトと、横軸 Hz 目盛り・縦軸 dB 目盛りの描画方法のレビュー
\end{itemize}

\section*{2. AIへの指示（プロンプト）}
AIへ送った指示のうち、主要なものを以下に示す。

\begin{lstlisting}
発展課題は「Processing でスペクトルアナライザを作り、マイクから入力された音信号の波形を
スペクトル解析した結果をウィンドウに描画する」プログラムを作りたい。
ベースにしたいファイル:
- Arduino:  work/shiozawa/work-0422/HCK02_02/HCK02_02.ino（動作確認済み、6.67 kHz サンプリング）
- Processing: work/shiozawa/work-0422/HCK02_02_processing/HCK02_02.pde（楽曲再生＋時間波形描画）

要件:
- Arduino 側はそのまま流用してよい（1 行 1 サンプル送信、150 µs 周期）
- Processing 側でシリアル受信したサンプルを蓄積し、minim の FFT クラスで解析する
- スペクトルはバーグラフ表示。横軸は Hz、縦軸は dB
- 楽曲再生のキーバインド（'p' / '1'-'6'）はそのまま残す
\end{lstlisting}

\begin{lstlisting}
FFT のサイズと窓関数、表示する周波数範囲を決めたい。
- サンプリングレート 6666.67 Hz、ナイキストは約 3333 Hz
- 周波数分解能と応答速度のバランスを考えて FFT_SIZE を決めたい
- 振幅の dB 変換時に log10(0) で発散しないようにしたい
\end{lstlisting}

\section*{3. AIの出力の使い方}
\begin{itemize}
  \item \cboxoff そのまま使った
  \item \cboxon 一部修正して使った
  \item \cboxoff 参考にしただけ
\end{itemize}
$\rightarrow$ 上で選んだ理由：

AI が挙げた設計のうち、以下は妥当と判断して採用した：
\begin{itemize}
  \item Arduino 側は \texttt{HCK02\_02.ino} と同じく 1 行 1 サンプル送信を維持し、
        Processing 側の FFT 用バッファに直接流し込む構成
  \item \texttt{FFT\_SIZE} = 1024（2 のべき乗）。サンプリング 6666.67\,Hz では
        1 回の FFT に必要な時間が約 154\,ms で、約 6.5\,Hz のビン幅となり、
        スペアナの応答性と分解能のバランスがとれる
  \item Hamming 窓を適用してスペクトル漏れを抑える
        （\texttt{fft.window(FFT.HAMMING)}）
  \item \texttt{20 * log10(max(magnitude, 1e-6))} の形で dB 変換し、
        ゼロ入力時の $-\infty$ 発散を防ぐ
  \item 上段 (時間領域) ／ 下段 (周波数領域) の 2 段レイアウト、
        下段は横軸 Hz・縦軸 dB の目盛りを併記
\end{itemize}
一方で、AI から追加提案された「FFT 結果を対数周波数軸で表示する」案や
「ピークホールドの減衰アニメーション」は、初版では機能を最小に絞りたかったため
今回は採用を見送り、後から拡張する余地として残した。

\section*{4. 正しいと確認した方法}
\todo{実機テスト後に具体的な観測内容を記入する}

想定している確認手順：
\begin{itemize}
  \item Arduino IDE で \texttt{HCK02\_03.ino} がコンパイルでき、
        シリアルモニタ (921600\,bps) に数値が 1 行ずつ流れることを確認する
  \item Processing IDE で \texttt{HCK02\_03.pde} を起動し、
        上段に時間波形、下段にスペクトルバーグラフが描画されることを観察する
  \item PC スピーカから \texttt{'p'} で楽曲を再生し、マイクへ入力したとき、
        スペクトルが楽曲の音高 (C4 $\approx$ 262\,Hz, G4 $\approx$ 392\,Hz, A4 = 440\,Hz)
        付近で立ち上がることを目視で確認する
  \item \texttt{'1'} (正弦波) で再生したときに、基音 1 本が突出した綺麗な
        スパイク状スペクトルになることを確認する
        $\rightarrow$ 倍音の少ない音色と FFT 結果が直感的に一致するかを見る
  \item \texttt{'3'} (のこぎり波) や \texttt{'4'} (矩形波) に切り替えると、
        基音に加えて倍音が高い周波数まで段々と並ぶことを確認する
        $\rightarrow$ 倍音構成の理論と実測の対応を確認する
  \item 拍手や口笛を入力し、無音時とのスペクトル差が見えることを確認する
  \item \texttt{SAMPLE\_RATE\_HZ}（Processing 側）と \texttt{SAMPLE\_US}（Arduino 側）の
        対応が \texttt{1000000.0 / SAMPLE\_US} で正しく一致していることを
        ソース上で計算しなおして確認する
\end{itemize}

\section*{5. 問題や間違い}
\begin{itemize}
  \item \cboxoff あった
  \item \cboxoff なかった
\end{itemize}
\todo{実機テスト後に、実際に気づいた問題や AI の誤り・自分が気づいた点を記入する}

検証予定ポイント：
\begin{itemize}
  \item シリアル受信が間に合わずに FFT バッファに穴が空くと、スペクトルが乱れる。
        \texttt{frameRate} と \texttt{port.available()} を表示し、取りこぼしが
        起きていないか確認する
  \item Arduino のサンプリングジッタが大きいと FFT の周波数軸が歪む。
        既知の音叉や PC で生成した既知周波数の正弦波（例: 440\,Hz）を入れ、
        スペクトルのピーク位置が想定ビン (440 / 6.51 $\approx$ 67 番目) に来るか確認する
  \item 起動直後に \texttt{fftInput} がゼロ初期化のまま FFT にかかるため、
        最初の数フレームは底値スペクトルが見える。
        想定どおり徐々にサンプルが埋まるかを観察する
  \item Processing のシリアルポート名 (\texttt{/dev/cu.usbmodem...}) は PC ごとに違うため、
        メンバ B の環境で書き換え忘れによる起動失敗が起きないか
\end{itemize}

\section*{6. 自分で工夫した点}
\begin{itemize}
  \item Arduino 側 (\texttt{HCK02\_03.ino}) は \texttt{HCK02\_02.ino} を継承して
        サンプリング周期 150\,µs を維持し、Processing 側の \texttt{SAMPLE\_RATE\_HZ}
        を \texttt{1000000.0 / SAMPLE\_US} で導出する形にして、両者の数値が
        ずれない構造にした
  \item Processing 側 (\texttt{HCK02\_03.pde}) では、ADC 値をシリアル受信した
        その場で \texttt{[-1, +1]} に正規化して直流成分を抜き、
        FFT 入力バッファ \texttt{fftInput} と時間波形バッファ \texttt{waveBuf} の
        両方に書き込む構成にした $\rightarrow$ 1 サンプルで時間領域・周波数領域の
        両方を更新できる
  \item FFT の縦軸を \texttt{20 * log10(max(magnitude, 1e-6))} で dB 化し、
        \texttt{SPECTRUM\_DB\_MIN} / \texttt{SPECTRUM\_DB\_MAX} の範囲外を
        \texttt{constrain} でクランプすることで、ゼロ入力時にバーが
        画面下に飛ぶ／対数の発散で描画が壊れる、といった事故を防いだ
  \item 表示する最高周波数 \texttt{DISPLAY\_MAX\_HZ} を 3000\,Hz に絞り、
        ナイキスト直近の使い物にならない領域を捨てて
        可聴域（楽器音域）に画面を集中させた
  \item バーの色を周波数比 \texttt{i / numBins} で連続的に変化させ、
        どの帯域が立っているかを色でも認識しやすくした
        （初期方針の「落ち着いた配色」の範囲内で、彩度を控えめに調整）
  \item 楽曲再生のキーバインド (\texttt{'p'} / \texttt{'1'}--\texttt{'6'}) を
        \texttt{HCK02\_02.pde} から維持し、波形 → 周波数の対応を
        音色を切り替えながら直感的に検証できるようにした
\end{itemize}

\section*{参考文献}
\begin{itemize}
  \item 生成AI利用申告書テンプレート（ハッカソン1 講義資料, \texttt{references/lectures/}）
  \item 提出課題2 のコード一式（本リポジトリ \texttt{work/shiozawa/work-0422/HCK02\_02/} 配下）
  \item 動作確認済み Arduino: \texttt{work/shiozawa/work-0422/HCK02\_02/HCK02\_02.ino}
  \item Processing Reference: \texttt{Serial}, \texttt{bufferUntil}, \texttt{serialEvent}
        \url{https://processing.org/reference/libraries/serial/}
  \item Minim Reference: \texttt{FFT}, \texttt{AudioOutput}, \texttt{Oscil}, \texttt{Line}
        \url{https://code.compartmental.net/minim/}
  \item Minim FFT Javadoc: \url{https://code.compartmental.net/minim/javadoc/ddf/minim/analysis/FFT.html}
\end{itemize}

\end{document}
